\documentclass[../../main.tex]{subfiles}
\begin{document}
\section{对称张量与交错张量}\label{section:ch12-对称与交错张量}

一般而言，重排协变张量的自变量，其值未必有可预测的变化。但某些特殊张量——例如点积——重排自变量时值不变；另一些—— notably 行列式——交换两个自变量时变号。本节描述两类在自变量重排时变化简单的张量：对称张量与交错张量。

\subsection*{对称张量}

设 $V$ 为有限维向量空间。称 $V$ 上的协变 $k$ 张量 $\alpha$ 为\textbf{对称的}，若交换任意一对自变量其值不变：
\begin{align*}
\alpha(v_1,\dots,v_i,\dots,v_j,\dots,v_k)=\alpha(v_1,\dots,v_j,\dots,v_i,\dots,v_k)
\end{align*}
对所有 $1\leqslant i<j\leqslant k$。

\begin{exercise}\label{exercise:ch12-13}
证明对协变 $k$ 张量 $\alpha$，下列条件等价：
\begin{enumerate}[(a)]
\item $\alpha$ 对称；
\item 对任意向量 $v_1,\dots,v_k\in V$，$\alpha(v_1,\dots,v_k)$ 在 $v_1,\dots,v_k$ 任意重排下不变；
\item 在任意基下，$\alpha$ 的分量 $\alpha_{i_1\dots i_k}$ 在指标任意置换下不变。
\end{enumerate}
\end{exercise}

$V$ 上对称协变 $k$ 张量之集是 $T^k(V^*)$ 的线性子空间，记为 $\Sigma^k(V^*)$。

存在从 $T^k(V^*)$ 到 $\Sigma^k(V^*)$ 的自然投影，定义如下。先令 $S_k$ 为 $k$ 个元素的对称群，即集合 $\{1,\dots,k\}$ 的置换群。给定 $k$ 张量 $\alpha$ 与置换 $\sigma\in S_k$，定义新 $k$ 张量 $\sigma\alpha$ 为
\begin{align*}
(\sigma\alpha)(v_1,\dots,v_k)=\alpha(v_{\sigma(1)},\dots,v_{\sigma(k)}).
\end{align*}
注意 $\tau(\sigma\alpha)=(\tau\sigma)\alpha$，其中 $\tau\sigma$ 表示 $\tau$ 与 $\sigma$ 的复合。我们定义投影 $\mathrm{Sym}\colon T^k(V^*)\to\Sigma^k(V^*)$，称为\textbf{对称化}，为
\begin{align}
\mathrm{Sym}\,\alpha=\frac{1}{k!}\sum_{\sigma\in S_k}\sigma\alpha.\label{align:ch12-对称化定义}
\end{align}
更明确地，即
\begin{align*}
(\mathrm{Sym}\,\alpha)(v_1,\dots,v_k)=\frac{1}{k!}\sum_{\sigma\in S_k}\alpha(v_{\sigma(1)},\dots,v_{\sigma(k)}).
\end{align*}

\begin{proposition}[对称化的性质]\label{proposition:ch12-对称化性质}
设 $\alpha$ 是有限维向量空间上的协变张量。
\begin{enumerate}[(a)]
\item $\mathrm{Sym}\,\alpha$ 是对称的；
\item $\mathrm{Sym}\,\alpha=\alpha$ 当且仅当 $\alpha$ 对称。
\end{enumerate}
\end{proposition}

\begin{proof}
设 $\alpha\in T^k(V^*)$。若 $\tau\in S_k$ 是任意置换，则
\begin{align*}
(\mathrm{Sym}\,\alpha)(v_{\tau(1)},\dots,v_{\tau(k)})
=\frac{1}{k!}\sum_{\sigma\in S_k}(\sigma\alpha)(v_{\tau(1)},\dots,v_{\tau(k)})
=\frac{1}{k!}\sum_{\sigma\in S_k}\alpha(v_{\tau\sigma(1)},\dots,v_{\tau\sigma(k)}).
\end{align*}
令 $\rho=\tau\sigma$，当 $\sigma$ 遍历 $S_k$ 时 $\rho$ 也遍历 $S_k$，故上式等于
\begin{align*}
\frac{1}{k!}\sum_{\rho\in S_k}\alpha(v_{\rho(1)},\dots,v_{\rho(k)})=(\mathrm{Sym}\,\alpha)(v_1,\dots,v_k).
\end{align*}
这表明 $\mathrm{Sym}\,\alpha$ 对称。

若 $\alpha$ 对称，则由习题~\ref{exercise:ch12-13}(b)，对每个 $\sigma\in S_k$ 有 $\sigma\alpha=\alpha$，故立即得 $\mathrm{Sym}\,\alpha=\alpha$。反之，若 $\mathrm{Sym}\,\alpha=\alpha$，则由 (a) $\mathrm{Sym}\,\alpha$ 对称，故 $\alpha$ 对称。
\end{proof}

若 $\alpha,\beta$ 是 $V$ 上的对称张量，$\alpha\otimes\beta$ 一般不对称。但用对称化算子可定义一种新积，它把一对对称张量变为另一个对称张量。若 $\alpha\in\Sigma^k(V^*)$，$\beta\in\Sigma^l(V^*)$，定义它们的\textbf{对称积}为 $(k+l)$ 张量 $\alpha\beta$（两个张量并列、中间无积号），由
\begin{align*}
\alpha\beta=\mathrm{Sym}(\alpha\otimes\beta)
\end{align*}
给出。更明确地，$\alpha\beta$ 对向量 $v_1,\dots,v_{k+l}$ 的作用为
\begin{align*}
\alpha\beta(v_1,\dots,v_{k+l})=\frac{1}{(k+l)!}\sum_{\sigma\in S_{k+l}}\alpha(v_{\sigma(1)},\dots,v_{\sigma(k)})\beta(v_{\sigma(k+1)},\dots,v_{\sigma(k+l)}).
\end{align*}

\begin{proposition}[对称积的性质]\label{proposition:ch12-对称积性质}
\begin{enumerate}[(a)]
\item 对称积对称且双线性：对所有对称张量 $\alpha,\beta,\gamma$ 与所有 $a,b\in\mathbb{R}$，
\begin{align*}
\alpha\beta=\beta\alpha,\qquad (a\alpha+b\beta)\gamma=a\alpha\gamma+b\beta\gamma=\gamma(a\alpha+b\beta).
\end{align*}
\item 若 $\alpha,\beta$ 是余向量，则
\begin{align*}
\alpha\beta=\frac{1}{2}(\alpha\otimes\beta+\beta\otimes\alpha).
\end{align*}
\end{enumerate}
\end{proposition}

\begin{exercise}\label{exercise:ch12-16}
证明命题~\ref{proposition:ch12-对称积性质}。
\end{exercise}

\subsection*{交错张量}

继续设 $V$ 为有限维实向量空间。称 $V$ 上协变 $k$ 张量 $\alpha$ 为\textbf{交错的}（亦称反对称、斜对称），若交换任意两个自变量时其值变号。即对所有 $v_1,\dots,v_k\in V$ 与每对不同指标 $i,j$，有
\begin{align*}
\alpha(v_1,\dots,v_i,\dots,v_j,\dots,v_k)=-\alpha(v_1,\dots,v_j,\dots,v_i,\dots,v_k).
\end{align*}
交错协变 $k$ 张量也称为\textbf{外形式}、\textbf{多余向量}或 $k$ 余向量。$V$ 上所有交错协变 $k$ 张量构成的子空间记为 $\Lambda^k(V^*)\subseteq T^k(V^*)$。

回忆对任意置换 $\sigma\in S_k$，$\sigma$ 的\textbf{符号} $\mathrm{sgn}\,\sigma$ 当 $\sigma$ 为偶置换（可表为偶数个对换之复合）时为 $+1$，奇置换时为 $-1$。下面的习题是习题~\ref{exercise:ch12-13}~的类比。

\begin{exercise}\label{exercise:ch12-17}
证明对协变 $k$ 张量 $\alpha$，下列条件等价：
\begin{enumerate}[(a)]
\item $\alpha$ 交错；
\item 对任意向量 $v_1,\dots,v_k$ 与任意置换 $\sigma\in S_k$，
\begin{align*}
\alpha(v_{\sigma(1)},\dots,v_{\sigma(k)})=(\mathrm{sgn}\,\sigma)\alpha(v_1,\dots,v_k);
\end{align*}
\item 在任意基下，$\alpha$ 的分量 $\alpha_{i_1\dots i_k}$ 交换任意两个指标时变号。
\end{enumerate}
\end{exercise}

每个 $0$ 张量（即一个实数）既对称又交错，因为没有自变量可交换。类似地，每个 $1$ 张量既对称又交错。$V$ 上交错 $2$ 张量就是斜对称双线性型。值得注意：每个协变 $2$ 张量 $\beta$ 都可表为一个交错张量与一个对称张量之和，因为
\begin{align*}
\beta(v,w)=\frac{1}{2}\big(\beta(v,w)-\beta(w,v)\big)+\frac{1}{2}\big(\beta(v,w)+\beta(w,v)\big)=\alpha(v,w)+\eta(v,w),
\end{align*}
其中 $\alpha(v,w)=\frac{1}{2}(\beta(v,w)-\beta(w,v))$ 是交错张量，$\eta(v,w)=\frac{1}{2}(\beta(v,w)+\beta(w,v))$ 是对称张量。对更高阶张量这不再成立，如习题~12-7~所示。

交错张量也有对称化与对称积的类似物，但我们推迟到第~14~章再引入，届时将更详细研究交错张量的性质。
\end{document}
